\documentclass[UTF8,a4paper,11pt,openany]{ctexbook}
\usepackage{../common/statlearnbook}
\SLSetBookMetadata{第 5 册}{统计学习方法讲义 v2.7.0}{FIG-P609-01 SA2 局部页验证}
\begin{document}
\setcounter{page}{644}
\setcounter{chapter}{32}
\setcounter{figure}{8}
\pagestyle{headings}
\chaptermark{马尔可夫链蒙特卡罗法}
\sectionmark{收敛与样本诊断}
有限样本的方差等效有效样本量只在分母为正时定义。用截断经验自相关替换真实自相关只能得到诊断估计，不能单独证明收敛。

图\ref{fig:V5-C03-autocorrelation-ess}把经验ACF与同一预设窗口下的有限样本加权ESS并列为一个诊断单元。
\input{../../绘图源码/第05册_采样方法主题模型与图排序/V5-C03/fig_v5_c03_autocorrelation_ess.tex}
\FloatBarrier
读图时，先看左面板在 $k=1,\ldots,6$ 的预设窗口内保留了哪些正经验自相关；再看右面板如何用有限样本权重形成 $\widehat\tau_{6,n}$；最终得到正相关增大方差权重、并使同长度轨迹的 $\widehat N_{\mathrm{eff}}$ 减小。

\paragraph{诊断边界。}有限轨迹的经验ACF和ESS都是诊断量，必须与核的可达性、返回性、非周期性及所需矩条件分开核验。
\end{document}
